% !TEX program = xelatex
\documentclass[UTF8]{ctexart}

\usepackage{amsmath,amssymb}
\usepackage[a4paper,margin=18mm]{geometry}
\usepackage{enumitem}
\usepackage{fancyhdr}
\usepackage{lastpage}

% TikZ for the schematic figure in Q11
\usepackage{tikz}
\usetikzlibrary{calc,angles,quotes}

\begin{document}
	\section{圆锥曲线的常用技巧}
	\begin{enumerate}
		\item 齐次化
		\item 齐次化升级
		\item 斜率双用
		\item 定比点差
		\item 极点极线
		\item 仿射变换
		\item 硬解定理
		\item 调和线束
		\item 点差法
		\item 蝴蝶模型
		\item 非对称韦达
		\item 焦点三角形
		\item 第三定义与点差法
		\item 蒙日圆、米勒圆
		\item 焦半径公式
		\item 阿基米德三角形
		
	\end{enumerate}


	\section*{阿基米德三角形（Archimedes Triangle）知识点总结（以抛物线为主）}
	
	\subsection*{1.\ 基本定义与记号}
	\textbf{抛物线上的阿基米德三角形：}
	在同一条抛物线 $\mathcal P$ 上取两点 $A,B$，分别作过 $A,B$ 的切线，相交于点 $P$。
	由两条切线段 $PA,PB$ 与弦 $AB$ 围成的三角形 $\triangle PAB$ 称为\textbf{阿基米德三角形}，其中 $AB$ 常称为其\textbf{底边}。
	
	\textbf{一般圆锥曲线的表述（补充）：}
	对任意非退化圆锥曲线（椭圆/双曲线/抛物线），若 $A,B$ 为切点、$P$ 为两切线交点，则由 $PA,PB$ 与 $AB$ 围成的三角形同样常被称为“阿基米德三角形”（很多题目实质是“切点弦/极线”问题，见第6节）。
	
	\subsection*{2.\ 阿基米德引理（几何核心结构）}
	在阿基米德三角形 $\triangle PAB$ 中，设 $M$ 为底边 $AB$ 的中点。
	过 $P$ 作中线 $PM$，令它与抛物线再交于点 $O$（通常 $O$ 位于弓形内部）。
	过 $O$ 作抛物线切线，分别交 $PA,PB$ 于 $A_1,B_1$。
	
	则有（常称“阿基米德引理”）：
	\begin{enumerate}
		\item $PM \parallel$ 抛物线的轴；
		\item $A_1,B_1$ 分别是 $PA,PB$ 的中点（即 $PA_1=A_1A,\ PB_1=B_1B$）；
		\item $O$ 是 $PM$ 的中点（即 $PO=OM$），从而 $A_1B_1$ 是 $\triangle PAB$ 的中位线，且 $A_1B_1 \parallel AB$。
	\end{enumerate}
	
	\subsection*{3.\ 面积结论：抛物线把阿基米德三角形按 $2:1$ 分割}
	设由弦 $AB$ 与抛物线围成的弓形面积为 $S_{\text{seg}}$，阿基米德三角形面积为 $S_{\triangle PAB}$。
	经典结论为：
	\[
	S_{\text{seg}}=\frac{2}{3}\,S_{\triangle PAB},
	\quad\text{等价于}\quad
	S_{\triangle PAB}-S_{\text{seg}}=\frac{1}{3}\,S_{\triangle PAB}.
	\]
	也常表述为：\textbf{抛物线在阿基米德三角形内部截出的部分与外侧剩余部分面积比为 $2:1$。}
	
	\textbf{与“抛物线求积(Quadrature of the Parabola)”的关系：}
	阿基米德证明了：\textbf{弓形面积等于某个内接三角形面积的 $\frac{4}{3}$ 倍}。
	而“阿基米德三角形”是他进行穷竭法分割时最自然的外包三角形之一，因此在很多教材/竞赛题里，$2/3$ 与 $4/3$ 往往会配套出现。
	
	\subsection*{4.\ 解析几何常用模板（以 $y^2=2px$ 为例）}
	取标准抛物线
	\[
	\mathcal P:\ y^2=2px\quad(p>0),
	\]
	其焦点 $F\left(\frac p2,0\right)$，准线 $x=-\frac p2$，轴为 $x$ 轴。
	
	\subsubsection*{4.1\ 切线方程与两切线交点坐标}
	若 $A(x_1,y_1)\in\mathcal P$，则 $x_1=\dfrac{y_1^2}{2p}$，切线方程为
	\[
	\ell_A:\ yy_1=p(x+x_1).
	\]
	同理，$B(x_2,y_2)\in\mathcal P$ 的切线
	\[
	\ell_B:\ yy_2=p(x+x_2).
	\]
	令 $\ell_A\cap\ell_B=P(x_P,y_P)$，可解得
	\[
	y_P=\frac{y_1+y_2}{2},\qquad
	x_P=\frac{y_1y_2}{2p}.
	\]
	
	\subsubsection*{4.2\ “中线平行于轴”的代数秒证}
	底边中点 $M\!\left(\dfrac{x_1+x_2}{2},\dfrac{y_1+y_2}{2}\right)$，
	而 $y_P=\dfrac{y_1+y_2}{2}$，故 $P,M$ 纵坐标相同，
	从而 $PM$ 为水平线，\textbf{平行于 $x$ 轴}，即平行于抛物线轴。
	
	\subsubsection*{4.3\ 参数法（高效解题补充）}
	令参数 $t\in\mathbb R$，抛物线可参数化为
	\[
	A(t):\ \left(\frac p2 t^2,\ pt\right),
	\quad\text{切线：}\quad
	ty=x+\frac p2 t^2.
	\]
	若 $A=A(t_1),\,B=B(t_2)$，则两切线交点
	\[
	P=\ell_{t_1}\cap \ell_{t_2}:
	\quad
	x_P=\frac p2 t_1t_2,\qquad
	y_P=\frac p2 (t_1+t_2).
	\]
	这个模板在“求交点/求轨迹/证垂直”类题中非常常用。
	
	\subsection*{5.\ 底边过焦点的特殊情形（“焦点阿基米德三角形”）}
	仍以 $y^2=2px$ 为例，设弦 $AB$ \textbf{过焦点} $F$，并作阿基米德三角形 $\triangle PAB$。
	则常用的结论链为：
	\[
	AB\ \text{过}\ F
	\quad\Longleftrightarrow\quad
	t_1t_2=-1
	\quad\Longrightarrow\quad
	\begin{cases}
		\text{(i) } \ell_{t_1}\perp \ell_{t_2}\ \ (\triangle PAB\ \text{在}\ P\ \text{处直角}),\\[2pt]
		\text{(ii) } x_P=\dfrac p2 t_1t_2=-\dfrac p2\ \ (\Rightarrow P\ \text{在准线上}),\\[2pt]
		\text{(iii) } PF\perp AB.
	\end{cases}
	\]
	
	其中 (iii) 可用参数直接算斜率验证：
	\[
	k_{AB}=\frac{pt_2-pt_1}{\frac p2(t_2^2-t_1^2)}=\frac{2}{t_1+t_2},\qquad
	k_{PF}=\frac{0-\frac p2(t_1+t_2)}{\frac p2-(-\frac p2)}=-\frac{t_1+t_2}{2},
	\]
	故 $k_{AB}\cdot k_{PF}=-1$，即 $PF\perp AB$。
	
	\subsection*{6.\ 进一步补充：它本质上是“切点弦/极点极线”}
	对任意圆锥曲线 $\mathcal C$：
	若从外点 $P$ 向 $\mathcal C$ 作两条切线，切于 $A,B$，
	则\textbf{连线 $AB$ 称为 $P$ 关于 $\mathcal C$ 的极线（polar），点 $P$ 称为直线 $AB$ 的极点（pole）}。
	这意味着：
	\begin{enumerate}
		\item 许多“阿基米德三角形”的轨迹题，本质是“极线随点动而绕某定点转动”的射影性质；
		\item 若一个点在某条线的极线上，则该线的极点在该点的极线上（La Hire 定理的互逆性）。
	\end{enumerate}
	因此：\textbf{遇到“切线交点、切点弦、过焦点/准线、直角”等关键词时，往往可以尝试极点极线视角简化。}
	
	\subsection*{7.\ 常见高考/竞赛题型与套路（补充）}
	\begin{enumerate}
		\item \textbf{证明平行/中点：} 优先用第4节交点公式，直接比较中点坐标；
		\item \textbf{直角与准线：} 一旦出现“弦过焦点”，立刻尝试参数条件 $t_1t_2=-1$，可一行得到 $P$ 在准线、两切线垂直；
		\item \textbf{面积：} 看到“弓形面积”优先联想到 $4/3$（内接三角形）与 $2/3$（阿基米德三角形）两条比例；
		\item \textbf{轨迹：} 交点 $P$ 的坐标经常能写成 $t_1+t_2,\ t_1t_2$，再用题设约束把它们消去得到轨迹方程。
	\end{enumerate}
	
	\subsection*{参考资料（含用户指定链接）}
	\begin{thebibliography}{99}
		\bibitem{ctk-arch} Cut-the-Knot: \textit{Archimedes Triangle and Squaring of Parabola}.
		\bibitem{ctk-lambert} Cut-the-Knot: \textit{Circumcircle of Three Parabola Tangents (Lambert's theorem)}.
		\bibitem{wiki-quad} 维基百科：\textit{抛物线求积（Quadrature of the Parabola）}.
		\bibitem{geogebra-focal} GeoGebra: \textit{Parabola: Tangents to the Endpoints of a Focal Chord}.
		\bibitem{cnblogs-focal} 博客园：\textit{阿基米德三角形（焦点弦情形的解析推导）}.
		\bibitem{pole-polar} Wikipedia: \textit{Pole and polar}.
		%		\bibitem{zhihu-user} （用户指定）\url{https://zhuanlan.zhihu.com/p/405295545}.
		%		\bibitem{weixin-user} （用户指定）\url{https://mp.weixin.qq.com/s/55Gwh23a5vOKqdstkkL7ag}.
	\end{thebibliography}
	
	
\end{document}